\section{Three Edits}
\label{sec:ch06-three-edits}

\index{HotChocolate.ApolloFederation@\code{HotChocolate.ApolloFederation}!the package|(}%
One package reference goes into \code{Sessions.csproj}, on the same version as
the rest of the family:

\begin{minted}[highlightlines={12},fontsize=\footnotesize]{xml}
<Project Sdk="Microsoft.NET.Sdk.Web">

  <PropertyGroup>
    <TargetFramework>net10.0</TargetFramework>
    <Nullable>enable</Nullable>
    <ImplicitUsings>enable</ImplicitUsings>
    <GenerateDocumentationFile>true</GenerateDocumentationFile>
    <NoWarn>$(NoWarn);CS1591</NoWarn>
  </PropertyGroup>

  <ItemGroup>
    <PackageReference Include="HotChocolate.ApolloFederation" Version="16.6.1" />
    <PackageReference Include="HotChocolate.AspNetCore" Version="16.6.1" />
    <PackageReference Include="HotChocolate.AspNetCore.CommandLine" Version="16.6.1" />
    <PackageReference Include="HotChocolate.Data.EntityFramework" Version="16.6.1" />
    <PackageReference Include="HotChocolate.Types.Analyzers" Version="16.6.1" />
    <PackageReference Include="Microsoft.EntityFrameworkCore.Sqlite" Version="10.0.11" />
  </ItemGroup>

</Project>
\end{minted}

\index{AddApolloFederation@\code{AddApolloFederation}}%
One call goes on the end of the builder chain in \code{Program.cs}. The line
above it gives up its semicolon and nothing else in the file moves:

\begin{minted}[highlightlines={34}]{csharp}
using ChilliCream.Nitro.App;
using HotChocolate.AspNetCore;
using Microsoft.EntityFrameworkCore;
using Sessions;

var builder = WebApplication.CreateBuilder(args);

// EF Core reports every command through the logging pipeline as well,
// at Information and with a duration attached. StatementLog below is
// this service's instrument, so silence the other one rather than have
// every statement printed twice.
builder.Logging.AddFilter(
    "Microsoft.EntityFrameworkCore.Database.Command",
    LogLevel.Warning);

builder.Services.AddDbContextFactory<ConferenceContext>(options =>
{
    options.UseSqlite("Data Source=conference.db");

    if (builder.Environment.IsDevelopment())
    {
        options.AddInterceptors(new StatementLog());
    }
});

builder
    .AddGraphQL()
    .AddSessionsTypes()
    .RegisterDbContextFactory<ConferenceContext>()
    .AddQueryContext()
    .AddPagingArguments()
    .AddGlobalObjectIdentification()
    .AddMutationConventions()
    .AddApolloFederation();

var app = builder.Build();

using (var scope = app.Services.CreateScope())
{
    var factory = scope.ServiceProvider
        .GetRequiredService<IDbContextFactory<ConferenceContext>>();
    using var db = factory.CreateDbContext();
    db.Database.EnsureCreated();
}

app.MapGraphQL()
    .WithOptions(options =>
    {
        // Serve the Nitro build that shipped in the package. The
        // default, ServeMode.Latest, downloads the current one
        // from ChilliCream's CDN instead.
        options.Tool.ServeMode = ServeMode.Embedded;
    });

app.RunWithGraphQLCommands(args);
\end{minted}

\index{key@\code{[Key]}!goes on the type class}%
\index{ReferenceResolver@\code{[ReferenceResolver]}!goes on the type class}%
Then two attributes on each type that is to become an entity. Both go on the
type class, the one carrying \code{[ObjectType<T>]}, and neither goes on a
method. \code{[Key]} names the field or fields that identify one of these,
written as
chapter~\ref{ch:federation-model} described a \code{FieldSet}.
\code{[ReferenceResolver]} names the method that will turn such a key back into
an object. Section~\ref{sec:ch06-the-key-arrives-as-a-string} prints both type
classes with both attributes on them, and writes the method the second one
names, which is where the chapter's work actually is.

\subsection{What the Schema Became}

\index{schema!what the package adds}%
Write the missing method, build, and export. Chapter~\ref{ch:federation-model}
said a subgraph is your service with four additions and named them; here they
are, as the difference between this exported schema and the one
chapter~\ref{ch:schema-design} left. Everything below was written by the
package. Nothing was typed into a schema file, because this book has no schema
file. As in chapters~\ref{ch:data-without-n-plus-1}
and~\ref{ch:schema-design}, the descriptions and \code{@cost} directives the
export writes are dropped here, and the fields of the two entity types are
elided, because what changed about them is on the first line:

\begin{graphqlsdl}
schema
  @link(
    url: "https://specs.apollo.dev/federation/v2.6"
    import: ["@key", "@shareable", "@tag", "FieldSet"]
  ) {
  query: Query
  mutation: Mutation
}

type Query {
  ...
  _service: _Service!
  _entities(representations: [_Any!]!): [_Entity]!
}

type _Service {
  sdl: String!
}

type Session implements Node @key(fields: "id") { ... }
type Speaker implements Node @key(fields: "id") { ... }

union _Entity = Session | Speaker

scalar _Any
scalar FieldSet
\end{graphqlsdl}

\index{link@\code{"@link}!the version is not chosen}%
The \code{@link} line is the opt-in. The url ends in a version, as
chapter~\ref{ch:federation-model} said it would, and that version is a default
rather than a decision: the package emits the one it emits unless you go looking
for the setting, and appendix~\ref{app:toolchain} is where the number lives. The
import list beside it holds the names the document actually uses, which is why
\code{@key} appears there the moment a type carries one.

\index{Entity@\code{\_Entity}!is generated from the keys}%
\code{\_Entity} is the union chapter~\ref{ch:federation-model} promised, and it
is generated rather than declared: two types carry a key, so two types are in
it, and a third key would put a third type in it. That is worth knowing because
it makes the union a readout. If a type you expected to be an entity is missing
from it, the key did not take, and the schema is telling you so in a line nobody
wrote.

\index{directives!the definitions travel with the schema}%
Below that, and elided here for length, the export also carries the definitions
of \code{@key}, \code{@link} and \code{@shareable}, written out in full with
their locations and their default arguments. A composer that reads this document
and nothing else knows what every directive in it means, which is
section~\ref{sec:ch06-the-schema-as-a-string}'s subject.
\index{HotChocolate.ApolloFederation@\code{HotChocolate.ApolloFederation}!the package|)}%
